\documentclass[11pt]{article}
\usepackage[T1]{fontenc}
\usepackage{textcomp}
\usepackage[margin=0.75in]{geometry}
\usepackage{amsmath,amssymb,amsthm}
\usepackage{array,booktabs}
\usepackage{hyperref}
\setlength{\parindent}{0pt}
\newtheorem{theorem}{Theorem}
\theoremstyle{definition}
\newtheorem{definition}{Definition}
\title{Flow Proof Artifact: prop-10-double-angle-triangle}
\author{Generated by \texttt{flow doc proof}}
\date{Flow Proof Book}
\begin{document}
\maketitle
To construct an isosceles triangle having each of the angles at the base double the remaining one.\\[0.5em]
\textbf{Source.} euclid — Elements, Book IV, Proposition 10\\[0.5em]
\bigskip
\noindent{\large \textbf{Derived fact 1.} \textit{To construct an isosceles triangle having each of the angles at the base double the remaining one} \hfill the Euclidean plane \cdot Euclid Book IV \cdot Proposition 10: to construct an isosceles triangle with base angles double the vertex}\par
\smallskip\noindent\textit{Coordinate: the Euclidean plane \cdot Euclid Book IV \cdot Proposition 10: to construct an isosceles triangle with base angles double the vertex}\par
\smallskip\noindent\textit{Source: euclid — Elements, Book IV, Proposition 10}\par
\smallskip\noindent\textit{Built on: proposition 4: side-angle-side congruence, for Book I of the Elements in on the Euclidean plane, proposition 32: the angle between tangent and chord equals the angle in the alternate segment, for Euclid Book III on the Euclidean plane}\par
\medskip
\noindent\textbf{Goal.} To construct an isosceles triangle having each of the angles at the base double the remaining one\par
\noindent$\displaystyle \text{isosceles triangle with double base angles}$\par
\medskip
\noindent
\begin{tabular}{@{} >{\raggedright\arraybackslash}p{0.06\textwidth} >{\raggedright\arraybackslash}p{0.47\textwidth} >{\raggedleft\arraybackslash}p{0.41\textwidth} @{}}
\textbf{\#} & \textbf{Proof} & \textbf{Mathematics} \\
\midrule
\textbf{1.} & We prove that proposition 10: to construct an isosceles triangle with base angles double the vertex for Euclid Book IV the Euclidean plane. & \\[0.45em]
\textbf{2.} & Let ABC = isosceles\_triangle. & \\[0.45em]
\textbf{3.} & Let angle\_A = vertex\_angle. & \\[0.45em]
\textbf{4.} & Let angle\_B = base\_angle. & \\[0.45em]
\textbf{5.} & From step 2, step 3, and step 4, we can deduce that angle B equals 2 times angle A. & $\displaystyle \angle B = 2 * \angle A$ \\[0.45em]
\textbf{6.} & We invoke the axiom governing Book I of the Elements in the Euclidean plane: proposition 4: side-angle-side congruence, for Book I of the Elements in on the Euclidean plane. & \\[0.45em]
\textbf{7.} & We invoke the derived fact governing Euclid Book III the Euclidean plane: proposition 32: the angle between tangent and chord equals the angle in the alternate segment, for Euclid Book III on the Euclidean plane. & \\[0.45em]
\textbf{8.} & From step 6 and step 7, this implies isosceles triangle with double base angles. Hence proven. & $\displaystyle \text{isosceles triangle with double base angles}$ \\[0.45em]
\bottomrule
\end{tabular}
\label{thm:Geometry-Euclid Book IV-proposition_10_to_construct_an_isosceles_triangle_with_base_angles_double_the_vertex}
\par\medskip\hrule
\end{document}
